\documentclass[10pt,a4paper]{amsart}
\usepackage[margin=19mm]{geometry}
\usepackage{amsmath,amssymb,mathtools}
\usepackage[scr=rsfs,cal=euler]{mathalfa}
\usepackage{xfrac}
\usepackage[unicode,hidelinks,bookmarks=false]{hyperref}
\newcommand{\ie}{i.e.\ }
\newcommand{\cf}{cf.\ }
\newcommand{\resp}{respectively\ }
\newcommand{\Subsection}{\subsection}
\providecommand{\nobreakdash}{\nobreak\hbox{-}}
\setlength{\emergencystretch}{2em}
\multlinegap0pt
\usepackage{amssymb}
\usepackage{braket}
\usepackage{mathtools}
\usepackage[shortlabels]{enumitem}
\newtheorem{lem}{Lemma}
\newcommand{\Gm}{\mathbb{G}_{m}}
\newcommand{\Z}{\mathbb{Z}}
\title{Galois Symbols for a Jacobian and Multiplicative Groups}
\author{Toshiro Hiranouchi, Rin Sugiyama}
\date{Source: arXiv:2608.11614v2 (CC BY 4.0)}
\begin{document}
\maketitle

\noindent\emph{Source and licence.} Galois Symbols for a Jacobian and Multiplicative Groups. Version arXiv:2608.11614v2 (CC BY 4.0), distributed by arXiv under the Creative Commons Attribution 4.0 International licence, http://creativecommons.org/licenses/by/4.0/, as recorded in the arXiv licence metadata for this exact version and shown on the abstract page https://arxiv.org/abs/2608.11614v2. Full licence notice: \texttt{licenses/arxiv-2608.11614-CC-BY-4.0.txt}.\par\smallskip

\noindent\textit{Editorial scope.} Counted unit: the article's lem lem:decomp\_CH (source0.tex lines 352--359). The article's complete original proof of it is delivered with it (source0.tex lines 360--479, one \texttt{proof} environment of 120 lines). No mathematics is written, repaired or re-derived: the statement and the proof are the article's own lines, in the article's own notation, and no retained line is substituted or re-flowed.  Retained uncounted as the source-local context the proof needs: lines 251--256; lines 268--271; lines 288--293.  Nothing was omitted from any retained span, so the package carries no omission record.  No retained line contains a citation, so the wrapper carries no bibliography and no bibliography entry.  No retained line contains a figure environment, an image-inclusion command or drawing code, so no image is delivered and the package declares no asset dependency.  Editorial formatting only, in addition: an AMS \texttt{amsart} preamble that restates the article's own declarations and macros used by the retained lines, namely the environments \texttt{lem} on separate counters, together with the standard packages those lines need.  The retained environments print in this document as Lemma 1 in order of appearance, whereas the article prints the same items with the numbering of its own full document.  The complete native source of the article as distributed in the arXiv e-print for this exact version is bundled unchanged as \texttt{sources/207432/source0.tex}.
\begin{equation}\label{eq:Akh}
 \varphi_r\colon
 CH^{r+1}(C,r)
 \overset{\sim}{\to}
 K_r\bigl(k;CH_0(C),\Gm\bigr).
\end{equation}

\begin{equation}\label{eq:Akh-sym}
 \varphi_r^{-1}\bigl(\set{z,a_1,\ldots,a_r}_{E/k}\bigr)
 = (p_E)_*\bigl(z\cap a_1\cap\cdots\cap a_r\bigr)
\end{equation}

\begin{equation}\label{eq:CH}
 CH^{r+1}(C,r)
 \simeq
 K_r^M(k)\oplus
 K_r(k;J,\Gm).
\end{equation}

\begin{lem}\label{lem:decomp_CH}
	For any $r\ge 0$, the isomorphism \eqref{eq:Akh} induces isomorphisms 
    \begin{align}
    &CH^{r+1}(h^0(C),r)=0,\\
    &CH^{r+1}(h^1(C),r)\simeq K_r(k;J,\Gm),\ \text{and}\\
    &CH^{r+1}(h^2(C),r)\simeq K_r^M(k).
    \end{align}
\end{lem}

\begin{proof}
%Since $h^0(C)\simeq (\operatorname{Spec}(k),1_k)$ and $h^2(C)\simeq\mathbb L$,
%we have
%\[
% CH^{r+1}(h^0(C),r)
% \simeq CH^{r+1}(k,r)=0
%\]
%and
%\[
% CH^{r+1}(h^2(C),r)
% \simeq CH^r(k,r)
% \simeq K_r^M(k).
%\]
%We identify the latter group with the Milnor $K$-theory summand in
%\eqref{eq:CH}. 
%
%Let $f\colon C\to\operatorname{Spec}(k)$  and $i\colon P\hookrightarrow C$ 
%be the structure morphism and the closed immersion, respectively.
%Let
%$\gamma\colon C\to C\times C$ given by $x\mapsto(x,P)$.
%Then
%\[
% \gamma_*[C]=[C\times P],\qquad
% p_1\circ\gamma=\operatorname{id}_C,\qquad
% p_2\circ\gamma=i\circ f.
%\]
%Hence, for $z\in CH^a(C,b)$, the projection formula gives
%\[
%\begin{aligned}
% (\pi_2)_*(z)
% &=(p_2)_*\bigl([C\times P]\cap p_1^*z\bigr)\\
% &=(p_2)_*\gamma_*\bigl(\gamma^*p_1^*z\bigr)\\
% &=(p_2\circ\gamma)_*(z)
% =i_*f_*(z).
%\end{aligned}
%\]
%Thus $(\pi_2)_*=i_*f_*$.
%
%Let $E/k$ be finite and let
%\[
% w=(p_E)_*
% \bigl(z\cap a_1\cap\cdots\cap a_r\bigr)
% \in CH^{r+1}(C,r),
%\]
%where $z\in CH_0(C_E)$ and $a_1,\ldots,a_r\in E^\times$.
%By the projection formula and the compatibility of proper
%push-forward with the Nesterenko--Suslin--Totaro isomorphism, we have
%\[
% f_*(w)
% =
% N_{E/k}\bigl(
% \deg_E(z)\{a_1,\ldots,a_r\}
% \bigr)
% \in K_r^M(k).
%\]
%Consequently,
%\[
% (\pi_2)_*(w)
% =
% i_*N_{E/k}\bigl(
% \deg_E(z)\{a_1,\ldots,a_r\}
% \bigr).
%\]
%Under Akhtar's isomorphism \eqref{eq:Akh}, this is the projection onto
%the summand $K_r^M(k)$ in \eqref{eq:CH}. Therefore, the kernel of
%$(\pi_2)_*$ corresponds to the summand $K_r(k;J,\Gm)$.
%
%Since the $h^0(C)$-part of $CH^{r+1}(C,r)$ is zero, the kernel of
%$(\pi_2)_*$ is precisely its $h^1(C)$-part. Hence
%\[
% CH^{r+1}(h^1(C),r)
% \simeq K_r(k;J,\Gm).
%\]
%
    Put $\alpha:=[P\times C], {}^t\alpha:=[C\times P]$.
    We first compute $\alpha_\ast, {}^t\alpha_\ast$ on $CH^1(C)$.
    For $z\in CH^1(C)$, we have
    $$
    \alpha\cap (z\times C)=0,\ {}^t\alpha \cap (z\times C)=z\times P
    $$
    which shows
    $$
    \alpha_\ast(z)=0,\ {}^t\alpha_\ast(z)=\deg(z)[P]\in \Z[P].
    $$
    This shows that the kernel of ${}^t\alpha_\ast$ on $CH^1(C)$ is equal to the kernel of the degree map
    $$
    \deg:CH^1(C) \to \Z
    $$ and proves the assertion in the case $r=0$.
    

    By \eqref{eq:Akh} and \eqref{eq:Akh-sym}, we take an element $w$ of $CH^{r+1}(C,r)$ of the following form
    $$
    w:=(p_E)_\ast\bigl(z\cap a_1\cap\cdots\cap a_r\bigr).
    $$
    By the base change formula for algebraic cycles with respect to proper push-forwards and flat pull-backs, it suffices to compute $\alpha_\ast(w)$ in case $E=k$.
    Then, we have
    \begin{align}
        \alpha_\ast(w)
        &={p_2}_\ast\bigl(\alpha\cap p_1^\ast(z\cap a_1\cap\cdots\cap a_r)\bigr)\\
        &={p_2}_\ast\bigl((\alpha\cap p_1^\ast z)\cap p_1^\ast(a_1\cap\cdots\cap a_r)\bigr)=0,
    \end{align}
    and
    \begin{align}
        {}^t\alpha_\ast(w)
        &={p_2}_\ast\bigl({}^t\alpha\cap p_1^\ast(z\cap a_1\cap\cdots\cap a_r)\bigr)\\
        &={p_2}_\ast\bigl(({}^t\alpha\cap p_1^\ast z)\cap p_1^\ast(a_1\cap\cdots\cap a_r)\bigr)\\
        &={p_2}_\ast\bigl([z\times P]\cap p_1^\ast(a_1\cap\cdots\cap a_r)\bigr)\\
        &=\deg(z)\bigl([P]\cap a_1\cap\cdots\cap a_r\bigr)\\
        &=\deg(z)\{a_1, \dots, a_r\}\in K_r^M(k).
    \end{align}
    Here in the last equality, we identify $CH^r(P,r)$ with $K^M_r(k)$.
    This shows that the kernel of ${}^t\alpha_\ast$ on $CH^{r+1}(C,r)$ is equal to the kernel of the map
    $$
    CH^{r+1}(C,r)\to K_r^M(k).
    $$
    Thus, the isomorphism \eqref{eq:Akh} induces
    $$
    CH^{r+1}(h^0(C),r)=0,\ CH^{r+1}(h^1(C),r)\simeq K_r(k;J,\Gm),\ CH^{r+1}(h^2(C),r)\simeq K_r^M(k).
    $$
\end{proof}

\end{document}
